\documentclass[11pt]{article}

\usepackage[margin=1in]{geometry}
\usepackage{graphicx}
\usepackage{booktabs}
\usepackage{amsmath}
\usepackage{amssymb}
\usepackage{hyperref}
\usepackage{caption}
\usepackage{subcaption}
\usepackage{xcolor}
\usepackage{placeins}
\usepackage{multirow}

\title{KABC Benchmark: Results (Experiments 1--4, Current Redesign)}
\author{Kernel-based Active Clustering of Bandit Arms\\
\small Thuot, Vogt, Ghoshdastidar, Verzelen (arXiv:2601.07535)}
\date{\today}

\begin{document}
\maketitle

\noindent This report summarizes the experimental setup, baseline algorithms,
and results for Experiments 1--4 of the \texttt{kernel-active-bandit}
benchmark suite. Experiments 1--3 stand as they were after a substantial
redesign of all three (algorithm roster, grids, and figures all changed from
the previous report). Full design rationale lives in
\texttt{docs/experiments.md}; this document summarizes what was actually run
and what the results show, plus a full pass at inferences. \textbf{All four
experiments are now complete at their full planned 50 reps.}
\textbf{Experiment 4 was substantially redesigned on 2026-09-16}
(Section~\ref{sec:setup-exp4}) and its full 50-rep production run under the
new design finished cleanly on 2026-09-17 --- see
Section~\ref{sec:results-exp4} for the real results, and the Run Provenance
paragraph below.

\paragraph{Run provenance.} \texttt{results/exp1/results.parquet} holds
6{,}300 rows (14 grid points $\times$ 9 algorithms $\times$ 50 reps) ---
\textbf{complete at the planned 50 reps}. 13 of these 14 points are the
canonical grid; the 14th ($N=90$) fills the gap between $N{=}60$ and
$N{=}120$ on the $N$ sweep, added permanently to \texttt{configs/exp1.yaml}.
\texttt{acb-gaussian} (Section~\ref{sec:acb-gaussian}) was added as a
ninth algorithm row on 2026-09-15, after the other eight's 50-rep run had
already completed --- merged in via a separate scratch-run-and-merge
(REVISION 6, \texttt{docs/experiments.md}), not a resume of the original
run.
\texttt{results/exp2/results.parquet} holds 2{,}400 rows (6 settings
$\times$ 8 algorithms $\times$ 50 reps) --- \textbf{complete at the planned
50 reps}. \texttt{results/exp3\_\{1,2,3\}/results.parquet} hold 300 rows each
(6 algorithms $\times$ 50 reps) --- \textbf{complete at the planned 50 reps}.
All three configs share the same underlying architecture: dedicated scripts
(\texttt{scripts/run\_experiment1/2/3.py}), calibrated threshold constants
(Section~\ref{sec:calibration}), and figures regenerated by
\texttt{notebooks/01\_gaussian\_mixtures.ipynb},
\texttt{notebooks/02\_same\_mean\_shapes.ipynb}, and
\texttt{notebooks/03\_movielens\_random\_pairs.ipynb} respectively. All three
experiments' \texttt{-0.9x} self-comparison rows were renamed to
\texttt{-0.5x} this round (Section~\ref{sec:baselines} explains why).
Experiment 4's redesigned \texttt{results/exp4\_baseline/results.parquet}
and \texttt{results/exp4\_amplified\_imbalance/results.parquet} each hold
100 rows (2 algorithms $\times$ 50 reps) --- \textbf{complete at the
planned 50 reps}, produced by the new \texttt{scripts/run\_experiment4.py};
figures regenerated by \texttt{notebooks/04\_movielens\_prototypes.ipynb}.
See Section~\ref{sec:results-exp4} for the real results.

\tableofcontents

\section{Experimental Setup}

\subsection{Experiment 1 --- Gaussian mixtures: competitiveness in the parametric regime}

\paragraph{Claim tested.} In the mean-separated Gaussian setting (home turf of
mean-based methods like ACB), KABC's budget is within a modest factor of ACB's;
budget scales as $s_*^{-2}\log(N/\delta)$ (the paper's Theorem 3.1). Also
(current redesign): how conservative each adaptive method's own stopping
rule is, and whether active/adaptive sampling beats naive uniform sampling
at a matched budget.

\paragraph{Instance family.} \texttt{GaussianMixtureEnv}: $N$ arms in $K=3$ balanced
clusters, arm $i$ in cluster $k$ distributed as $\mathcal{N}(\mu_k, \sigma^2 I_d)$,
$\sigma=1$. Cluster means $\mu_k = \frac{\Delta}{\sqrt2}(e_{2k-1}+e_{2k})$, giving each
pair of cluster means an actual Euclidean separation of exactly $\Delta$.

\paragraph{Grid.} One knob swept at a time, others held at the base config
$N{=}60, K{=}3, d{=}10, \Delta{=}3$:
\begin{itemize}
  \item $\Delta \in \{1, 1.5, 2, 3, 4, 6\}$ (6 points --- the headline sweep)
  \item $d \in \{6, 10, 50, 100\}$ (4 points; $d=2$ excluded, construction requires $d \ge 2K=6$)
  \item $N \in \{30, 60, 90, 120\}$ (4 points; $N=90$ added this round to fill
    the gap between 60 and 120)
\end{itemize}
14 grid points total (the base config is each sweep's own midpoint, so it
is run identically under all three sweep tags --- a documented, accepted
redundancy).

\paragraph{Algorithms: 8 rows per (grid point, rep).} Three independent BASE
algorithms --- \texttt{kabc}, \texttt{kabc-sg}, \texttt{acb} --- each run
once at the grid point's full \texttt{budget\_cap}. Each base is then,
separately, re-run at $\text{budget\_cap} = \min(0.5 \times \tau_{\text{own}},
\text{budget\_cap})$ with a fresh independent RNG substream ---
\texttt{kabc-0.5x}, \texttt{kabc-sg-0.5x}, \texttt{acb-0.5x} respectively,
each testing how conservative THAT algorithm's own stopping rule is.
\texttt{unif-kkm-1x} and \texttt{unif-kkm-0.5x} are tied specifically to
\texttt{kabc}'s own $\tau$ (not kabc-sg's or acb's), testing whether
active/adaptive selection beats naive uniform sampling at a matched budget.
No \texttt{tas-fw} here: its finite-alphabet assumption does not hold for
continuous Gaussian instances (Experiment~3 only).

\paragraph{Budget\_cap.} One value per grid point, shared across all 8
algorithm rows and every rep: $\text{budget\_cap} = 200\, N\, s_*^{-2}
\log(N/\delta)$ (Theorem~3.1), $s_*^2$ Monte-Carlo-estimated
(\texttt{kabc.metrics.estimate\_s\_star\_sq}).

\subsection{Experiment 2 --- Identical means, different shapes: nonparametric separation}
\label{sec:setup-exp2}

\paragraph{Claim tested.} When clusters share the same mean, mean-based
methods are blind (ARI at chance for any budget) while KABC recovers the
partition with budget governed by the MMD gap. Also (current redesign): the
same stopping-rule-conservativeness and active-vs-uniform-sampling questions
as Experiment 1.

\paragraph{Instance families and the current 6-setting grid.} Three
families, $N=40$, $K=2$ throughout. As of this report, the grid is
\textbf{6 fixed, illustrative settings} (not a sweep whose in-between
variation is the point --- see \texttt{docs/experiments.md}'s Experiment 2
GRID TRIM NOTE for the rationale):
\begin{itemize}
  \item \textbf{\texttt{scale\_pair}} (2a): cluster A $=\mathcal N(0,1)$, cluster B
    $=\mathcal N(0,\sigma_B^2)$ --- same mean, different \emph{variance}.
    $\sigma_B \in \{2, 3\}$.
  \item \textbf{\texttt{moment\_matched\_mixture}} (2b): cluster A $=\mathcal
    N(0,1)$, cluster B $= 0.5\,\mathcal N(-a,b^2) + 0.5\,\mathcal N(a,b^2)$,
    $b^2=1-a^2$ --- same mean \emph{and} variance, different \emph{shape}.
    $a \in \{0.85, 0.9\}$.
  \item \textbf{\texttt{ring\_blob}} (2c): cluster A $=$ a 2D ring (Gaussian
    noise $\sigma=0.1$), cluster B $=\mathcal N(0,I_2)$, both centered at the
    origin --- same mean, different \emph{2D geometry}. Ring radius $r \in
    \{2, 3\}$.
\end{itemize}
50 repetitions per setting (full planned run, see Run Provenance above).

\paragraph{Algorithms: 8 rows per (setting, rep), identical scheme to
Experiment 1.} Same base/0.5x-self-comparison/unif-kkm-tied-to-kabc
structure as Experiment 1 (see above). No \texttt{kabc-oracle} (dropped ---
its one-shot required budget exceeded \texttt{budget\_cap} at every grid
point of the old design, so the ``price of adaptivity'' figure it existed
for never produced real data). No \texttt{tas-fw}, same reason as
Experiment 1.

\paragraph{Budget\_cap.} Identical mechanism to Experiment 1: one
$s_*$-formula value per setting, shared across all 8 rows and every rep.

\subsection{Experiment 3 --- MovieLens random pairs: robustness across the real-data pair distribution}

\paragraph{Claim tested.} Whether KABC's advantage (and ACB's/TaS-FW's
relative behavior) holds up when the movie PAIR is randomly resampled every
repetition across the real diversity of the MovieLens catalog, under three
distinct pairing regimes that vary how much mean separation is present and
whether the second movie is cherry-picked (max MMD$^2$) or sampled
uniformly. Also: how conservative KABC's own stopping rule is.

\paragraph{Data.} \texttt{data/ml-1m-207/} --- the 207 MovieLens-1M movies
with $\ge 1000$ ratings. Discrete rating support $\{1,\ldots,5\}$,
RBF-on-ratings kernel, fixed median-heuristic bandwidth ($=1.0$) computed
once over the full 207-movie pool.

\paragraph{Instance.} $N=8$ arms, $K=2$ clusters, 4 arms per prototype
(movie), \texttt{MovieLensPairEnv}. For a given rep, a single shared
permutation of all 207 movie indices is drawn; each of the three versions
independently scans that same permutation and picks the first candidate
movie A with at least one eligible B under \emph{its own} gap rule (so
versions usually, but not always, land on the same A --- see
\texttt{docs/experiments.md} for the precise mechanism).

\paragraph{Version-specific B-selection rules.}
\begin{description}
  \item[\texttt{exp3\_1} --- tight mean gap, MMD-maximizing.] Candidates: all
    movies with $|\text{mean}_A - \text{mean}_B| \le 0.25$. $B = \arg\max
    \mathrm{MMD}^2$ among candidates.
  \item[\texttt{exp3\_2} --- moderate-to-large mean gap, uniform.]
    Candidates: all movies with $0.25 \le |\text{mean}_A - \text{mean}_B|
    \le 1.0$. $B$ sampled uniformly at random among candidates.
  \item[\texttt{exp3\_3} --- tight mean gap, uniform.] Candidates: all movies
    with $|\text{mean}_A - \text{mean}_B| \le 0.25$ (same tolerance as
    \texttt{exp3\_1}, a clean matched-tolerance ablation of cherry-picking
    vs.\ uniform B-selection). $B$ sampled uniformly at random.
\end{description}

\paragraph{Budget.} Fixed $\text{budget\_cap}=15{,}000$ pulls (not the
$s_*$-formula type --- not meaningful at this instance scale). Deliberately
not raised to clear \texttt{kabc}'s doubling-trick round-7 threshold
($28{,}560$ at $N=8,\delta=0.05$): this leaves \texttt{kabc}'s own
non-convergent behavior unchanged while giving \texttt{acb}/\texttt{tas-fw}
less runway.

\paragraph{Algorithms: 6 rows per (rep, version).} \texttt{kabc},
\texttt{kabc-0.5x}, \texttt{acb}, \texttt{tas-fw}, \texttt{unif-kkm-1x},
\texttt{unif-kkm-0.5x} --- the same base/0.5x/unif-kkm-tied-to-kabc scheme
as Experiments 1--2, restricted to the three algorithms for which it is
meaningful here (\texttt{kabc-sg} dropped this round --- its fragility on
this exact real-data setting is already thoroughly documented from a prior
run, see \texttt{docs/experiments.md}'s REVISION 4 note). 50 repetitions per
version (full planned run).

\subsection{Experiment 4 --- MovieLens prototypes: kabc's decision rule vs.\ unif-kkm's k-means under real imbalance}
\label{sec:setup-exp4}

\paragraph{Clean redesign (2026-09-16), superseding an earlier
4-condition/6-algorithm version entirely.} That version's own results are
not reported here; see \texttt{docs/experiments.md}'s Experiment 4 section
(git history) for what it found and exactly why each change below was made.

\paragraph{Claim tested.} Is kabc's decision rule (pairwise variance-aware
threshold + connected components) more robust than unif-kkm's k-means under
heterogeneous difficulty and real cluster-size imbalance? Both methods run
at the SAME (matched) budget on real MovieLens-1M rating distributions as
prototypes.

\paragraph{Two algorithms only.} \texttt{kabc}; \texttt{unif-kkm-1x}
(budget = that rep's own realized $\tau_{\texttt{kabc}}$, the matched-budget
design) --- the earlier version's $\tau_{\texttt{kabc}}$-blind fixed-budget
rows served their purpose (establishing unif-kkm loses even with MORE
budget than kabc used) but are not needed for this report.

\paragraph{Two settings, both real-data mean-gap-resampled every rep (no
hand-picked fixed instance).} Both mirror Experiment 3's own
\texttt{exp3\_3} mechanism (bounded mean-gap eligibility, uniform pick among
eligible, not MMD$^2$ cherry-picking): a close pair (mean gap $\le 0.25$)
and a far pair (mean gap $\ge 0.5$) guaranteed each rep, plus unconstrained
generic anchors. \texttt{baseline}: K=7, \texttt{far\_large}{=}12,
\texttt{far\_small}{=}2, \texttt{close\_a}{=}5, \texttt{close\_b}{=}4, 3
generic anchors (8, 6, 3) --- $N{=}40$. \texttt{amplified\_imbalance}: K=12,
\texttt{far\_large}{=}20, \texttt{far\_small}{=}2 (a 10:1 ratio, not
\texttt{baseline}'s 6:1), same close pair sizes, 8 generic anchors (10, 8,
7, 6, 5, 4, 3, 2) --- $N{=}76$; this is the regime where the earlier
design's K=12 condition showed the cleanest imbalance-specific signal
(unif-kkm's far-pair ARI below its close-pair ARI, small-prototype
misclassification $\sim$1.6--2$\times$ large-prototype), pushed further.

\paragraph{Memory guard override, real-verified.} Both settings run at
$n_{\text{jobs}}{=}5$ (not Experiments 1--3's 10--11) with
\texttt{kabc.fit}'s own additive, opt-in memory-guard override
(\texttt{memory\_guard\_n\_workers}{=}5,
\texttt{memory\_guard\_safety\_factor}{=}2 --- Experiments 1--3 are
entirely unaffected, keeping the original 11-worker/4$\times$ assumption).
This raises the safe per-iteration ceiling from $n_k{\sim}2852$ to
$n_k{\sim}5983$, verified with real measurements before use, not just the
relaxed formula: a real k=8-scale computation ($n{\approx}4602$, all pairs,
$N{=}40$) measured $627.7$ MB actual peak RSS per worker (vs.\ the pure
3-Gram-matrix model's $484.9$ MB); 5 real concurrent OS processes at that
scale measured $626.4$--$626.5$ MB each (sum $3.06$ GiB, wall-clock
$1242.4$s), with \texttt{vm\_stat} confirming zero new swapouts during the
run. This raises kabc's honest convergence ceiling from k$\le$7 to k$\le$8
for both settings ($budget\_cap$ raised correspondingly: $400{,}000$ for
\texttt{baseline}, $800{,}000$ for \texttt{amplified\_imbalance}, each
covering k=8's full cumulative cost with margin); k=9
($n_k{\approx}9200$--$9900$) was checked and rejected --- real projected
memory at 5 workers is $\sim$14.5 GiB, too close to this machine's 16GB
target. Full derivation in \texttt{docs/experiments.md}.

\paragraph{Metrics.} Standard $\tau$/\texttt{stopped}/ARI, plus a per-pair
ARI breakdown (\texttt{ari\_close\_pair}, \texttt{ari\_far\_pair}) and a
per-prototype-size misclassification rate
(\texttt{misclass\_rate\_large}/\texttt{\_small}, via best-overlap
alignment, median-split per instance) --- both carried over unchanged from
the earlier design.

\FloatBarrier
\section{Baseline Algorithms}
\label{sec:baselines}

\subsection{KABC and KABC-SG (this paper's own methods)}

KABC (Algorithm~2 of \cite{kabc2026}) runs a doubling-trick outer loop: at
iteration $k$, pull $n_k = \lceil 2^k \log(8(N^2-N)/\delta_k)\rceil$ fresh samples
per arm ($\delta_k = \delta/(4k^2)$), then build a same-cluster graph over all arm
pairs via a pairwise MMD-based test --- an edge $\{i,j\}$ is added iff the
empirical MMD statistic between arms $i$ and $j$ falls under a threshold. The
recovered partition is the graph's connected components; the loop stops the
first time the number of connected components equals the true $K$.

Two threshold variants are compared (identical outer loop, different pairwise
test):
\begin{description}
  \item[\texttt{kabc}] variance-aware threshold (Eq.~6 of \cite{kabc2026},
    building on the variance-aware kernel mean embedding estimates of
    \cite{wolferalquier2025}): sizes each pair's own threshold from that
    pair's \emph{empirical} RKHS variance, so a noisier pair automatically
    gets more slack.
  \item[\texttt{kabc-sg}] kernel-uniform sub-Gaussian threshold (the
    Discussion section's ablation in \cite{kabc2026}, substituting the
    minimax kernel-mean-embedding bound of \cite{tolstikhin2017},
    Proposition~A.1): uses only the fixed kernel bound (RBF kernel sup
    $=1$), identical for every pair regardless of that pair's own spread ---
    the ablation whose practical fragility Section~\ref{sec:conclusions}
    discusses in detail.
\end{description}

\paragraph{Implementation verification.} \texttt{src/kabc/thresholds.py}'s
\texttt{variance\_aware\_threshold} and \texttt{subgaussian\_threshold}
were checked directly against Eq.~6 and the Discussion section's
substitute formula in \cite{kabc2026} and match exactly, term for term
(including the $32/3$ coefficient in Eq.~6's second term and the
$(\sqrt{\cdot}+2)$ form of the sub-Gaussian bound) --- including correctly
distinguishing $\tilde g$ (the kernel's \emph{range}, $\sup-\inf$, used in
Eq.~6) from $\bar g$ (the kernel's \emph{supremum}, used in the
sub-Gaussian formula): these two quantities coincide numerically for the
RBF kernel used throughout this repo ($\tilde g=\bar g=1$, since the RBF
kernel's infimum is $0$ in the limit), but are conceptually distinct
quantities in the paper, and the code keeps them as separate named
parameters (\texttt{g\_tilde} vs.\ \texttt{g\_bar}) rather than silently
conflating them.

\paragraph{The \texttt{-0.5x} self-comparison rows (current redesign, all
three experiments).} For each base algorithm (\texttt{kabc},
\texttt{kabc-sg}, \texttt{acb}), a second, independent run at
$\text{budget\_cap} = \min(0.5\times\tau_{\text{own}}, \text{budget\_cap})$
tests how much margin that algorithm's own stopping rule carries. \textbf{Why
0.5x, not 0.9x:} an earlier version of this design used a 0.9x multiplier;
direct per-rep inspection of the resulting \texttt{kabc}/\texttt{kabc-0.9x}
data showed KABC's doubling-trick schedule quantizes its own realized $\tau$
to a small set of round-boundary checkpoints, and the FINAL round alone
typically costs roughly half of the total realized $\tau$. A 0.9x budget cut
therefore almost always landed short of that final round regardless of how
close to 1.0 the multiplier nominally was --- there was no way to construct
a real ``kabc re-run with 10\% less budget'' against this schedule; 0.9x and,
say, 0.6x tested essentially the same thing. 0.5x makes what is actually
being tested explicit: a meaningfully smaller, still-real budget cut. This
replaces an earlier, cruder $\text{budget\_multiplier}\in\{1.0,0.25\}$
restricted-budget pass applied uniformly across all algorithms.

\subsection{ACB --- Active Clustering with Bandit feedback}

Source: \cite{acb2024}. ACB* (this repo's adaptive variant, ``$\theta_*$ and
$\Delta_*$ unknown,'' referred to throughout as ACB) is a \emph{mean}-based active
clustering procedure (SRI + ADC subroutines): it estimates each arm's mean
directly (no kernel/MMD machinery) and clusters by Euclidean distance between
empirical means, adaptively refining which arms to pull based on estimated
separation. It is the ``home turf'' competitor for Experiment~1 (mean-separated
Gaussians, exactly ACB's design setting) and the mean-blindness foil for
Experiments~2--3 (same-mean, different-shape data, where a mean-only method has no
signal by construction).

\subsection{ACB-Gaussian --- the paper's own Gaussian-specialized, non-adaptive variant}
\label{sec:acb-gaussian}

Source: \cite{acb2024}, Appendix A. Beyond the main-text ACB* above, the
paper's own numerical experiments (its Figure~1, reporting $\le 1\%$ error
at $\delta=0.1$) do not actually use ACB*'s adaptive Algorithm~2 --- they
use a separate, non-adaptive variant (Algorithm~1, with $\Delta_*$ and
$\theta_*$ assumed known), with every concentration-inequality-derived
constant in Eq.~9--13 replaced by an exact $\delta$-level quantile: $n_{max}$
from a $\chi^2$ quantile, and $I$/$J$ from a normal quantile $\beta$ and a
Monte-Carlo-estimated product-quantile $\alpha$. We implement this variant
separately (\texttt{acb-gaussian}, \texttt{src/kabc/algorithms/
acb\_gaussian.py}) as a ninth, independent row in Experiment~1's algorithm
roster only (no \texttt{-0.5x} self-comparison counterpart).

\textbf{This baseline is evaluated with $\Delta$ and $\theta$ known,
matching the non-adaptive Algorithm~1 setting the paper's own authors use
for their Figure~1 experiment --- not the fully adaptive ACB* --- since
this is the fairest, most literal reproduction of the comparison the
paper's own authors ran.} We state this explicitly here rather than leaving
it implicit: a reader could otherwise reasonably wonder why one baseline
gets privileged information (the true separation $\Delta$ and balancedness
$\theta$) that no other algorithm in this benchmark receives. The answer is
that Experiment~1's \texttt{GaussianMixtureEnv} instances are
\emph{constructed from} these exact values (\texttt{configs/exp1.yaml}
passes $\Delta$ and $\sigma$ directly to the environment's constructor), so
nothing is being estimated from data that a fair comparison would otherwise
require estimating --- $\Delta$ and $\theta$ are known at instance-
construction time in the same sense $N$, $K$, and $d$ already are to every
algorithm in this repo. This does mean \texttt{acb-gaussian} does not match
CLAUDE.md's standard \texttt{fit(env, delta, budget\_cap, rng)} interface in
substance (only in call shape: $\Delta$/$\theta$/$\sigma$ are read directly
off \texttt{env}, not estimated), so it is wired only into Experiment~1's
own algorithm registry and hard-fails immediately if ever invoked on a
non-Gaussian environment, rather than silently misapplying a Gaussian-only
method outside its valid domain.

\textbf{No calibration.} Unlike \texttt{kabc}/\texttt{kabc-sg}/\texttt{acb},
\texttt{acb-gaussian} has no calibrated constant in
\texttt{configs/calibration.yaml} (the entry there is present and
explicitly documented as a deliberate, permanent absence, not a
not-yet-run placeholder): every quantity is an exact $\delta$-level
quantile of a known finite-sample distribution, not a concentration-
inequality bound with slack to tighten empirically.

\subsection{TaS-FW --- Track-and-Stop with Frank-Wolfe}

Source: \cite{tasfw2024}. Track-and-Stop based on Frank-Wolfe optimization, for
$N$-ary clustering of arms whose distributions live on a common \emph{finite}
alphabet. Only run in Experiment~3, where MovieLens ratings on
$\{1,\ldots,5\}$ satisfy the finite-alphabet assumption natively; not run on
Experiments~1--2's continuous instances, which would require an off-label
discretization approximation outside the algorithm's proven guarantees. Its
Track-and-Stop design re-scans its entire hypothesis space and re-solves a
linear program at \emph{every single pull} (not periodically, like KABC's
doubling trick), which is why Experiment~3 caps $N$ at 8: the hypothesis
space grows combinatorially in $N$, and empirically \texttt{tas-fw} already
costs two to three orders of magnitude more wall-clock per repetition than
every other algorithm at $N=8$ (confirmed again at the full 50 reps,
Section~\ref{sec:results-exp3}).

\subsection{Unif-KKM --- uniform sampling + kernel $k$-means (folklore baseline)}

No paper --- a standard, non-adaptive baseline implemented directly per the
design doc. Pulls a fixed number of samples per arm, then runs ordinary
kernel $k$-means on the resulting empirical kernel mean embeddings. Has no
confidence-based stopping rule and no $\delta$-PAC guarantee --- it always
spends its entire allotted budget and always ``completes''
($\text{stopped}=\text{True}$ by construction), so it is included purely as
a non-adaptive reference point, not as a $\delta$-correct competitor.
\texttt{unif-kkm} no longer sweeps an independent budget --- its two rows,
\texttt{unif-kkm-1x} and \texttt{unif-kkm-0.5x}, are tied specifically to
\texttt{kabc}'s own realized $\tau$ for that same rep
($\text{budget\_cap} = \min(1.0\text{ or }0.5 \times \tau_{\texttt{kabc}},
\text{budget\_cap})$), directly testing whether active/adaptive selection
beats naive uniform sampling at a matched budget.

\subsection{Related active-clustering methods not evaluated}
\label{sec:related-not-evaluated}

Several other recent papers address closely related active/bandit
clustering problems but were not included as baselines, either because
their problem setting differs from ours in a way that would make a
head-to-head comparison uninformative or unfair, or because of
prohibitive computational cost at our benchmark's scale.

\textbf{Optimal Clustering with Bandit Feedback} \cite{boc2024} proposes
BOC, an asymptotically optimal Track-and-Stop-style algorithm for the same
fixed-confidence clustering problem we study. We do not include it as a
baseline because Track-and-Stop methods require solving a constrained
optimization problem at every sampling round, making per-decision cost
prohibitive at our problem sizes; we observed comparable cost issues with
\texttt{tas-fw} \cite{tasfw2024}, which we retain only on a small reduced
instance (Experiment~3, $N=8$) for this reason.

\textbf{Clustering Items through Bandit Feedback: Finding the Right
Feature out of Many} \cite{banditclustering2025} studies a distinct
observation model in which the learner selects both an item and a
specific feature at each round, rather than observing a full feature
vector per pull. This partial-observation setting is not directly
comparable to our full-vector-feedback setting without a nontrivial
reformulation of the environment, and the method (BanditClustering) is
designed specifically around exploiting the sparsity of the discriminative
signal across features, which is not the structure our instances test.

\textbf{Near Optimal Best Arm Identification for Clustered Bandits}
\cite{clusteredbai2025} addresses a fundamentally different goal: given
$N$ agents each facing their own $K$-armed bandit problem, with agents
implicitly clustered by which bandit problem they share, the objective is
to identify each agent's best arm, not to recover the clustering as an end
in itself. Clustering is a means to accelerate best-arm identification via
information-sharing across agents in the same cluster, whereas partition
recovery is our sole objective; the two problems are not directly
comparable on the same metric.

\textbf{Online Clustering of Data Sequences with Bandit Information}
\cite{atboc2025} proposes ATBOC for a setting close to ours
(fixed-confidence clustering of $M$ arms into $K$ groups from bandit
feedback), permitting non-identical within-cluster distributions and
vector observations. We did not implement it as a baseline in the current
benchmark; a natural extension of this work would include it as an
additional parametric competitor alongside \texttt{acb}, since its
algorithmic structure (average tracking) differs from both ACB's
mean-based approach and BOC's Track-and-Stop approach.

\textbf{Efficient Clustering in Stochastic Bandits} \cite{ebc2026}
proposes EBC, which addresses the same per-round optimization bottleneck
that makes BOC and \texttt{tas-fw} expensive, by taking a single
gradient-like step toward the optimum at each round rather than solving
the full optimization problem, while retaining asymptotic optimality.
This is a promising, more efficient alternative to BOC that emerged
concurrently with our own work; we did not have the opportunity to
incorporate it into the present benchmark, but it is a natural inclusion
for future comparison given it targets the exact same computational
bottleneck we cite as our reason for excluding BOC and limiting
\texttt{tas-fw}'s scope.

\FloatBarrier
\section{Calibration}
\label{sec:calibration}

\texttt{kabc}, \texttt{kabc-sg}, and \texttt{acb} each have a theoretically-derived
but conservative leading constant in their stopping threshold. Per the
calibration protocol (\texttt{docs/experiments.md} \S0.2), each constant was
empirically calibrated once (binary search for the smallest constant keeping
empirical error $\le \delta/2$, validated under the REAL formula-derived
budget regime, not an artificially generous one), then frozen and reused
unchanged across every experiment. The frozen values actually used for this
report:

\begin{table}[h]
\centering
\begin{tabular}{lccc}
\toprule
Algorithm & $c_{\text{cal}}$ & Calibration-set error & Held-out error \\
\midrule
\texttt{kabc} & 0.103516 & 0.000 & 0.000 \\
\texttt{kabc-sg} & 0.121094 & 0.000 & 0.000 \\
\texttt{acb} & 0.658984 & 0.020 & 0.010 \\
\texttt{tas-fw} & --- (no calibratable constant) & --- & --- \\
\bottomrule
\end{tabular}
\end{table}

\paragraph{Important caveat.} Calibration instances are a \emph{single}
fixed Gaussian-mixture instance (\texttt{GaussianMixtureEnv}'s construction
is deterministic given $N,K,d,\Delta$ --- there is no way to draw ``20
different instances,'' only 20 independently-seeded noise draws of the one
instance). The frozen constants above were validated against that one
instance's own real budget regime, \emph{not} against every real grid point
in Experiments~1--3. Section~\ref{sec:conclusions} shows this still matters
directly: \texttt{kabc-sg}'s calibration does not fully generalize to the
harder grid points/settings in either experiment below.

\subsection{Implementation robustness: two properties worth stating explicitly}
\label{sec:robustness}

\paragraph{The graceful trivial-partition fallback gives ARI exactly
$0.0$, never negative, by construction --- not noise.} When a
budget-exhausted run has insufficient evidence to split arms, \texttt{kabc}
and \texttt{acb} both fall back to the trivial ``everyone in one cluster''
answer (\texttt{stopped=False}) rather than guessing. ARI \emph{can} go
negative in general (a maximally scrambled alternating partition against a
balanced true partition gives ARI $<0$), but the adjusted-for-chance ARI
formula makes a degenerate single-cluster partition's observed and
expected agreement coincide exactly, so this specific fallback always
scores exactly $0.000$ against any true partition. This is why
Experiment~2's \texttt{acb}/\texttt{acb-0.5x} rows land on ARI \emph{exactly}
$0.000$ at every one of the 6 settings (Table in
Section~\ref{sec:results-exp2}) rather than some small positive residual:
same-mean data gives \texttt{acb} no signal, it exhausts its budget every
time, and the fallback's exact-zero property is what produces the clean
result, not an artifact of rounding or aggregation.

\paragraph{A degenerate-bandwidth edge case on real MovieLens data was
found and fixed with a graceful stop, not a crash.} During Experiment~3
development, \texttt{median\_heuristic\_bandwidth}'s pilot-phase call
raised on a genuinely pathological real movie pair (Schindler's List vs.\
Rear Window: both extremely skewed toward high ratings, where a 5-per-arm
pilot draw sampled only tied $\{3,4,5\}$ values by chance, making the
median pairwise distance exactly $0$). This is fixed at the source
(\texttt{src/kabc/kernels.py}, \texttt{src/kabc/algorithms/kabc.py}):
\texttt{median\_heuristic\_bandwidth} now accepts \texttt{on\_degenerate=
"mean"}, falling back to the mean pairwise distance (still informative
unless \emph{every} pair is tied) with a \texttt{UserWarning}; if the mean
is \emph{also} degenerate, \texttt{kabc} now stops gracefully
(\texttt{stopped=False}, \texttt{stop\_reason="degenerate\_pilot"}) rather
than crashing, matching the existing budget-cap/memory-guard graceful-stop
convention. This matters beyond a single pair: real rating data can be
arbitrarily skewed, and a method whose bandwidth estimation crashes on
such an instance --- rather than degrading gracefully --- would silently
lose an entire batch of in-flight repetitions in a parallel run (confirmed
directly against this exact case; \texttt{kabc} now converges correctly on
it instead of crashing the whole batch).

\FloatBarrier
\section{Results --- Experiment 1 (50 reps, 14 grid points, complete)}
\label{sec:results-exp1}

Aggregated over all 14 grid points, 50 reps each (700 runs/algorithm,
6{,}300 total rows, 9 algorithms):

\begin{table}[h]
\centering
\begin{tabular}{lccccc}
\toprule
Algorithm & Stopped rate & Error rate & Mean ARI & Mean $\tau$ \\
\midrule
\texttt{kabc}          & 100.0\% &  0.0\% & 1.000 &  9{,}646 \\
\texttt{unif-kkm-1x}   & 100.0\% &  0.0\% & 1.000 &  9{,}646 \\
\texttt{unif-kkm-0.5x} & 100.0\% &  0.0\% & 1.000 &  4{,}812 \\
\texttt{acb}           & 100.0\% & 10.4\% & 0.990 &  8{,}018 \\
\texttt{acb-gaussian}  & 100.0\% &  0.0\% & 1.000 &  1{,}152 \\
\texttt{kabc-sg}       &  67.1\% & 32.9\% & 0.881 & 79{,}853 \\
\texttt{kabc-sg-0.5x}  &  17.4\% & 82.6\% & 0.479 & 36{,}964 \\
\texttt{acb-0.5x}      &  11.0\% & 92.7\% & 0.106 &  3{,}731 \\
\texttt{kabc-0.5x}     &   3.0\% & 97.0\% & 0.047 &  3{,}799 \\
\bottomrule
\end{tabular}
\end{table}

\begin{figure}[h]
\centering
\includegraphics[width=0.7\textwidth]{../results/figures/exp1_fig1_tau_vs_delta.png}
\caption{Stopping budget $\tau$ vs.\ $\Delta$, \texttt{kabc}/\texttt{kabc-sg}/\texttt{acb}/
\texttt{acb-gaussian} only (the four independent base algorithms, each using the grid
point's own full, undivided \texttt{budget\_cap}), with the theoretical
$\tau \propto s_*^{-2}$ bound overlaid (dashed).}
\label{fig:exp1-fig1}
\end{figure}

\begin{figure}[h]
\centering
\begin{subfigure}{0.48\textwidth}
\includegraphics[width=\textwidth]{../results/figures/exp1_fig2_tau_vs_d.png}
\caption{$\tau$ vs.\ ambient dimension $d$.}
\end{subfigure}
\hfill
\begin{subfigure}{0.48\textwidth}
\includegraphics[width=\textwidth]{../results/figures/exp1_fig2b_tau_vs_N.png}
\caption{$\tau$ vs.\ number of arms $N$ (4 points).}
\end{subfigure}
\caption{Budget scaling along the two secondary sweep dimensions.}
\label{fig:exp1-fig2}
\end{figure}

\begin{figure}[h]
\centering
\begin{subfigure}{0.32\textwidth}
\includegraphics[width=\textwidth]{../results/figures/exp1_fig3_convergence_vs_delta.png}
\caption{vs.\ $\Delta$}
\end{subfigure}
\hfill
\begin{subfigure}{0.32\textwidth}
\includegraphics[width=\textwidth]{../results/figures/exp1_fig3b_convergence_vs_d.png}
\caption{vs.\ $d$}
\end{subfigure}
\hfill
\begin{subfigure}{0.32\textwidth}
\includegraphics[width=\textwidth]{../results/figures/exp1_fig3c_convergence_vs_N.png}
\caption{vs.\ $N$}
\end{subfigure}
\caption{Convergence rate (fraction of reps with $\text{stopped}=\text{True}$), \texttt{kabc}/\texttt{kabc-sg}/\texttt{acb}/\texttt{acb-gaussian} only.}
\label{fig:exp1-fig3}
\end{figure}

\begin{figure}[h]
\centering
\includegraphics[width=0.75\textwidth]{../results/figures/exp1_fig4_ari_vs_delta.png}
\caption{ARI vs.\ $\Delta$, all 9 algorithm rows.}
\label{fig:exp1-fig4}
\end{figure}

\begin{figure}[h]
\centering
\begin{subfigure}{0.48\textwidth}
\includegraphics[width=\textwidth]{../results/figures/exp1_fig4b_ari_vs_d.png}
\caption{ARI vs.\ $d$.}
\end{subfigure}
\hfill
\begin{subfigure}{0.48\textwidth}
\includegraphics[width=\textwidth]{../results/figures/exp1_fig4c_ari_vs_N.png}
\caption{ARI vs.\ $N$.}
\end{subfigure}
\caption{ARI along the two secondary sweep dimensions, all 9 algorithm rows.}
\label{fig:exp1-fig4bc}
\end{figure}

\paragraph{Where \texttt{kabc-sg} struggles --- confirmed at the full 50
reps.} At $\Delta=1$ (the hardest separation), \texttt{kabc-sg} converges
only $8\%$ of the time (4/50 reps stopped, up slightly from the pilot's
0/10) and its mean $\tau$ balloons to $272{,}173$ --- close to the
doubling trick's $277{,}440$ ceiling at this $N,\delta$ --- while
\texttt{kabc} converges every time at $\tau=30{,}170$. The pattern is even
starker at the two hardest $d$ points ($d=50,100$): \texttt{kabc-sg} never
converges at either (0/50 both, mean $\tau=277{,}440$ pinned exactly at
the ceiling both times), while \texttt{kabc} converges every time at a
small fraction of that budget ($14{,}100$ both points). The two points
diverge sharply in ARI despite both having $0\%$ stopped: $d=50$'s ARI is
$0.531$ (a partially-correct partition even without formal convergence)
while $d=100$'s is $0.033$ (close to the graceful trivial-partition
fallback) --- \texttt{kabc-sg} is not just failing to converge at $d=100$,
it is failing to find useful structure at all.

\paragraph{ACB's error is concentrated at $\Delta\le 2$, not spread evenly
--- a more precise claim than the aggregate figure.} The $10.4\%$ aggregate
error rate (table above) obscures a sharp pattern along the $\Delta$
sweep: per-point error rates are $18\%$ ($\Delta=1$), $10\%$ ($\Delta=1.5$),
$18\%$ ($\Delta=2$), $2\%$ ($\Delta=3$), $0\%$ ($\Delta=4$), $0\%$
($\Delta=6$). \texttt{acb} is essentially error-free once the mean gap
clears $\Delta=3$, and all of its error mass sits at the three smallest,
hardest gaps in the sweep --- consistent with a mean-based method losing
precision as cluster means move closer together, not a uniform
error rate that happens to average to $10$--$11\%$.

\paragraph{At $\Delta=1$, \texttt{kabc}'s budget-competitiveness with
\texttt{acb} lands exactly where \cite{kabc2026}'s own Section~4 analytic
special case predicts it should, not as an unexplained pattern.} Section~4
gives an explicit Gaussian/linear-kernel comparison: KABC's budget bound
matches the known-optimal Gaussian linear-kernel rate of \texttt{acb}'s
own paper (\cite{acb2024}) whenever $\|\mu_i-\mu_j\|\le\sigma_j^2/\sqrt{\bar
g}$ and the ambient dimension is moderate. Experiment~1's instances use
$\sigma=1$ throughout and the RBF kernel's supremum $\bar g=1$, so this
condition reduces to $\Delta\le 1$ --- exactly the smallest, hardest point
on the $\Delta$ sweep (base $d=10$, well within the grid's own
``moderate'' $d\in\{6,\ldots,100\}$ range), sitting right at the boundary
of the regime the theorem addresses rather than comfortably inside or
outside it. Consistent with this: at $\Delta=1$ (now at the full 50-rep
$\Delta$-sweep), \texttt{kabc}'s realized $\tau$ ($30{,}170$) is already
\emph{lower} than \texttt{acb}'s ($62{,}697$) on the identical instances
--- KABC is not merely ``within a modest factor'' of ACB here but ahead of
it, precisely in the one regime the paper's own analysis singles out as
where KABC's bound should be competitive with ACB's optimal rate. We do
not read this as a tight numerical match (the theorem concerns the two
algorithms' theoretical \emph{bounds}, not their realized empirical
$\tau$, and \texttt{acb}'s own calibrated threshold has its own
conservativeness independent of this comparison) --- only that the
qualitative pattern (KABC staying budget-competitive with, here even
ahead of, ACB specifically at the grid's smallest gap) lands where the
paper's own theory says it should, rather than being an unexplained
empirical coincidence.

\paragraph{$N=90$ fits cleanly between the existing $N=60$/$N=120$ points;
its apparent pilot-scale anomaly is now confirmed as noise, not a real
effect.} \texttt{kabc}'s $\tau$ at $N=90$ ($9{,}360$) sits between its
$N=60$ ($6{,}000$) and $N=120$ ($12{,}960$) values, consistent with the
$\propto N$ scaling in the budget formula. The pilot (10 reps) had flagged
\texttt{kabc-sg} at $N=90$ as ``notably better-behaved'' than at $N=60$ or
$N=120$; at the full 50 reps this distinction has vanished:
\texttt{kabc-sg} reaches $100\%$ stopped and ARI $1.000$ at \emph{all
three} of $N=90$, $N=120$ ($100\%$/$1.000$), and $N=60$ ($98\%$/$0.999$)
--- essentially indistinguishable, all near ceiling. The real story on
this sweep is instead $N=30$ (the smallest, not previously flagged):
\texttt{kabc-sg} converges only $78\%$ of the time there (ARI $0.974$),
the actual outlier on this sweep once all four points are measured at
full statistical power. This confirms the pilot's $N=90$ observation was
exactly what it was flagged as being at the time: pilot-scale sampling
noise, not a real effect worth further investigation.

\paragraph{The \texttt{-0.5x} rows now show a much cleaner, starker
margin-collapse signal than the earlier 0.9x design --- confirmed at the
full 50 reps.} Every base algorithm's ARI drops sharply at 50\% of its own
realized $\tau$: \texttt{kabc} $1.000\to0.047$, \texttt{kabc-sg}
$0.881\to0.479$, \texttt{acb} $0.990\to0.106$ --- all consistent with the
pilot's $0.027$/$0.458$/$0.076$. Unlike the old 0.9x design's confusing
bimodal behavior (documented in \texttt{docs/experiments.md}'s MULTIPLIER
CHANGED note), 0.5x consistently and unambiguously demonstrates that all
three stopping rules carry little slack at this scale.
\texttt{unif-kkm-0.5x} barely moves ($1.000\to1.000$ exactly, aggregate),
since it is tied to \texttt{kabc}'s $\tau$, not its own, and half of
\texttt{kabc}'s budget is still ample for non-adaptive kernel $k$-means at
these grid points --- see Section~\ref{sec:conclusions} for why this
divergence between \texttt{kabc-0.5x} and \texttt{unif-kkm-0.5x} is itself
an interesting finding, not a contradiction.

\paragraph{\texttt{acb-gaussian} (Section~\ref{sec:acb-gaussian}) is
perfect on every single one of the 700 runs, at a fraction of every other
algorithm's budget --- checked per grid point, not just in aggregate, the
same discipline the ACB per-$\Delta$ error breakdown (\S6.3) required.}
Across all 14 grid points and 50 reps ($N=700$), \texttt{acb-gaussian}
reaches $100\%$ stopped and ARI $1.000$ \emph{exactly}, on every
individual grid point, not just in aggregate --- matching \texttt{kabc}'s
own zero-error record (\S\ref{sec:conclusions}). Its budget advantage vs.\
\texttt{kabc} is a genuinely uniform pattern across the grid, not an
aggregate driven by a few points: the per-grid-point ratio
$\tau_{\texttt{kabc}}/\tau_{\texttt{acb-gaussian}}$ ranges $5.4\times$ to
$17.4\times$ (aggregate mean-$\tau$ ratio $8.4\times$ sits squarely inside
that range). The comparison against \texttt{kabc-sg} is a different
story, and stating a single ``$69\times$ cheaper'' aggregate would repeat
exactly the mistake the ACB per-$\Delta$ error breakdown (\S6.3) was added
to correct: the per-point ratio
$\tau_{\texttt{kabc-sg}}/\tau_{\texttt{acb-gaussian}}$ ranges $4.7\times$
to $343\times$, with the two extreme values ($343\times$ at $d=50$,
$295\times$ at $d=100$) driven entirely by \texttt{kabc-sg} pinning at the
doubling trick's $277{,}440$ ceiling there (\S6.2) --- excluding those two
points, the range is $4.7\times$--$79\times$, still wide but not
three-orders-of-magnitude wide. The aggregate $69\times$ figure is real
but should be read as ``\texttt{kabc-sg} is dramatically more expensive
specifically where it also fails to converge,'' not as a uniform
budget-efficiency gap. This is not a surprising result given what
\texttt{acb-gaussian} is: Algorithm~1 with $\Delta$ and $\theta$ known
exactly (see Section~\ref{sec:acb-gaussian}'s explicit framing of why this
comparison is fair here specifically) and every constant replaced by an
exact $\delta$-level quantile rather than a conservative concentration-
inequality bound --- it is the paper's own ``sharper constants'' claim
(Appendix A) confirmed directly, not merely reproduced qualitatively. The
budget gap versus \texttt{acb} itself (mean $\tau$ $8{,}018$, non-zero
$10.4\%$ error) is the cleanest illustration in this report of the
difference between the paper's proof-derived (conservative, error-prone at
this constant) main-text algorithm and its own exact-quantile,
oracle-parameter variant.

\FloatBarrier
\section{Results --- Experiment 2 (50 reps, 6 settings, complete)}
\label{sec:results-exp2}

\begin{figure}[h]
\centering
\includegraphics[width=\textwidth]{../results/figures/exp2_fig0_headline_settings.png}
\caption{The 6 settings' instance geometry, one panel each: \texttt{scale\_pair}
and \texttt{moment\_matched\_mixture} panels are exact analytic densities
(A vs.\ that setting's B); \texttt{ring\_blob} panels are real 2D sample
scatters from \texttt{RingBlobEnv} itself.}
\label{fig:exp2-fig0}
\end{figure}

Aggregated over all 6 settings, 50 reps each (300 runs/algorithm, 2{,}400
total rows):

\begin{table}[h]
\centering
\begin{tabular}{lccccc}
\toprule
Algorithm & Stopped rate & Error rate & Mean ARI & Mean $\tau$ \\
\midrule
\texttt{kabc}          & 100.0\% &  0.0\% & 1.000 &     33{,}430 \\
\texttt{unif-kkm-1x}   & 100.0\% &  7.3\% & 0.987 &     33{,}430 \\
\texttt{unif-kkm-0.5x} & 100.0\% &  9.7\% & 0.986 &     16{,}696 \\
\texttt{kabc-sg}       &  83.0\% & 47.7\% & 0.680 &     35{,}659 \\
\texttt{kabc-sg-0.5x}  &  42.0\% & 72.7\% & 0.495 &      8{,}309 \\
\texttt{kabc-0.5x}     &  12.3\% & 87.7\% & 0.123 &     15{,}721 \\
\texttt{acb}           &   0.0\% & 100.0\% & 0.000 & 1{,}326{,}538 \\
\texttt{acb-0.5x}      &   0.0\% & 100.0\% & 0.000 &   663{,}218 \\
\bottomrule
\end{tabular}
\end{table}

\begin{figure}[h]
\centering
\includegraphics[width=\textwidth]{../results/figures/exp2_fig1_budget_bars.png}
\caption{Stopping budget $\tau$, one bar plot, 6 blocks (settings) $\times$
4 bars (\texttt{kabc}, \texttt{kabc-sg}, \texttt{acb}, plus a
\texttt{theoretical bound} bar --- the setting's formula-based
\texttt{budget\_cap}, a single deterministic value with no error bar).}
\label{fig:exp2-fig1}
\end{figure}

\begin{figure}[h]
\centering
\includegraphics[width=\textwidth]{../results/figures/exp2_fig2_convergence_bars.png}
\caption{Convergence rate, 6 blocks $\times$ 3 bars (\texttt{kabc}/\texttt{kabc-sg}/\texttt{acb} only).}
\label{fig:exp2-fig2}
\end{figure}

\begin{figure}[h]
\centering
\includegraphics[width=\textwidth]{../results/figures/exp2_fig3_ari_bars.png}
\caption{ARI, 6 blocks $\times$ 8 bars (all algorithm rows).}
\label{fig:exp2-fig3}
\end{figure}

\paragraph{ACB is exactly at its theoretical budget bound everywhere, never
converging.} Figure~\ref{fig:exp2-fig1} shows \texttt{acb}'s bar coinciding
with the grey \texttt{theoretical bound} bar in every one of the 6 blocks
--- confirmed directly in the aggregate table ($\text{stopped}=0.0\%$,
$\text{ARI}=0.000$). All 6 settings are same-mean by construction, so a
mean-based method has no signal, spends its entire allotted budget every
time, and always falls back to the graceful trivial partition.

\paragraph{KABC-SG is family-dependent, and \texttt{moment\_matched\_mixture}
is uniformly the hardest family.} \texttt{kabc-sg} converges well on
\texttt{scale\_pair} and \texttt{ring\_blob} but its ARI collapses
specifically on \texttt{moment\_matched\_mixture} --- it \emph{converges}
confidently to a wrong-ish answer on this family far more than it fails to
converge, the same confident-wrong-answer failure mode documented for this
algorithm on real MovieLens data in the earlier (now-superseded) report.
\texttt{kabc}, by contrast, is at ARI $1.000$ on every one of the 6
settings individually, not just on average. This picture is confirmed, not
just reproduced, at the full 50 reps: \texttt{kabc-sg}'s aggregate ARI moved
from $0.616$ (10-rep pilot) to $0.680$ (50 reps) and its stopped rate from
$86.7\%$ to $83.0\%$ --- both well within pilot-scale noise, the same
qualitative pattern holds.

\paragraph{Precisely quantified: \texttt{kabc-sg}'s confidently-stopped
runs violate its own calibrated $\delta=0.05$ guarantee substantially on
\texttt{moment\_matched\_mixture} --- one of the strongest results in the
whole benchmark, not just an ablation weakness.} Restricting to
\texttt{kabc-sg}'s $95$ \texttt{stopped=True} runs on
\texttt{moment\_matched\_mixture} (out of 100 total reps across its two
settings), $79$ of them ($83.2\%$) have $\text{ARI}\neq 1$ --- a
confidently-wrong error rate over 16$\times$ higher than the $\delta=0.05$
guarantee \texttt{kabc-sg}'s calibrated threshold is supposed to certify.
\texttt{kabc} --- same family, same budget, same calibration protocol,
differing only in its variance-aware (rather than kernel-uniform)
pairwise threshold --- shows exactly $0$ errors on the same family. This
is direct empirical evidence for the variance-aware threshold's practical
value: it is not merely a modest accuracy edge, but the difference between
a stopping rule that actually holds its stated confidence level and one
that fails it by more than an order of magnitude on this family.
\cite{kabc2026}'s own Section~4 (``Benefit of variance-aware bounds'')
analytically anticipates exactly this failure mode: since the per-arm
RKHS variance $\mathcal V_i^*$ never exceeds the kernel's global bound
$\bar g$, the variance-aware bound is never worse than the sub-Gaussian
one and can substantially improve the budget whenever $\mathcal V_i^*\ll
\bar g$. \texttt{moment\_matched\_mixture} is constructed so that cluster
A and B share both mean \emph{and} variance and differ only in
higher-order shape (see Section~\ref{sec:setup-exp2}) --- exactly the
regime where the two per-arm RKHS variances used in Eq.~6 are informative
and small relative to $\bar g$, and where \texttt{kabc-sg}'s
kernel-uniform bound, blind to that per-arm variance information, pays
the largest possible price for ignoring it. The $83.2\%$ violation rate is
therefore not just an empirical curiosity turned up by this ablation: it
is the paper's own theoretically-predicted worst case for the sub-Gaussian
threshold, observed exactly where the theory says it should appear.

\paragraph{The \texttt{unif-kkm-0.5x} vs.\ \texttt{kabc-0.5x} gap
(Section~\ref{sec:conclusions} discusses the mechanism) is dramatic here
too, and sharper at 50 reps than the pilot suggested.}
\texttt{unif-kkm-0.5x}'s ARI ($0.986$) is essentially unchanged from
\texttt{unif-kkm-1x}'s ($0.987$), while \texttt{kabc}'s own ARI at 50\% of
its realized $\tau$ collapses to $0.123$. (The pilot's $0.067$ for
\texttt{kabc-0.5x} was itself within the range of a 10-rep estimate of this
50-rep $0.123$ --- the direction and scale of the effect were already
correct in the pilot.)

\FloatBarrier
\section{Results --- Experiment 3 (50 reps, 3 versions, complete)}
\label{sec:results-exp3}

\begin{figure}[h]
\centering
\includegraphics[width=\textwidth]{../results/figures/exp3_fig0_headline_pairs.png}
\caption{The real movie pair actually sampled at $\text{rep}=0$, for each of
the three versions, as real rating histograms (\texttt{data/ml-1m-207/ratings.parquet},
not recomputed from any algorithm output --- same category as Experiment~2's
headline figure, which plots real 2D sample scatters directly from \texttt{RingBlobEnv}).
At this rep, all three versions share the same shared-candidate-order first
eligible movie A (``Rain Man'', 1988) --- so this one rep directly shows how each
version's own $B$-selection rule (Section~\ref{sec:results-exp3}'s setup) picks a
different movie $B$ against the identical $A$: \texttt{exp3\_1}'s cherry-picked,
max-MMD$^2$ choice (Pulp Fiction, MMD$^2{=}0.052$) is visibly the most separated
of the three; \texttt{exp3\_3}'s uniformly-chosen $B$ under the same tight mean-gap
tolerance (Big, MMD$^2{=}0.009$) is visibly the most overlapping --- a direct,
real-data illustration of why \texttt{exp3\_3} is the hardest of the three versions
for every algorithm (\S\ref{sec:conclusions}).}
\label{fig:exp3-fig0}
\end{figure}

\begin{table}[h]
\centering
\caption{Experiment 3, $N=8$, 50 reps/algorithm/version. Stopped rate, mean $\tau$, mean ARI, mean wall-clock/rep.}
\begin{tabular}{llcccc}
\toprule
Version & Algorithm & Stopped rate & Mean $\tau$ & Mean ARI & Mean wall-clock/rep \\
\midrule
\multirow{6}{*}{\texttt{exp3\_1}}
 & \texttt{kabc}        &  100.0\% &  3{,}750 & 1.000 & 0.50\,s \\
 & \texttt{kabc-0.5x}   &   14.0\% &  1{,}719 & 0.140 & 0.25\,s \\
 & \texttt{acb}         &    0.0\% & 14{,}997 & 0.000 & 0.66\,s \\
 & \texttt{tas-fw}      &   38.0\% & 13{,}660 & 1.000 & 576.27\,s \\
 & \texttt{unif-kkm-1x} &  100.0\% &  3{,}750 & 1.000 & 0.20\,s \\
 & \texttt{unif-kkm-0.5x} & 100.0\% &  1{,}873 & 1.000 & 0.12\,s \\
\midrule
\multirow{6}{*}{\texttt{exp3\_2}}
 & \texttt{kabc}        &   98.0\% &  3{,}076 & 0.995 & 0.42\,s \\
 & \texttt{kabc-0.5x}   &   14.0\% &  1{,}397 & 0.140 & 0.21\,s \\
 & \texttt{acb}         &   20.0\% & 13{,}252 & 0.124 & 0.57\,s \\
 & \texttt{tas-fw}      &   70.0\% &  9{,}942 & 1.000 & 409.04\,s \\
 & \texttt{unif-kkm-1x} &  100.0\% &  3{,}076 & 1.000 & 0.16\,s \\
 & \texttt{unif-kkm-0.5x} & 100.0\% &  1{,}537 & 0.990 & 0.09\,s \\
\midrule
\multirow{6}{*}{\texttt{exp3\_3}}
 & \texttt{kabc}        &   80.0\% & 10{,}380 & 0.791 & 1.57\,s \\
 & \texttt{kabc-0.5x}   &   14.0\% &  4{,}942 & 0.135 & 0.58\,s \\
 & \texttt{acb}         &    0.0\% & 14{,}997 & 0.000 & 0.64\,s \\
 & \texttt{tas-fw}      &    4.0\% & 14{,}928 & 1.000 & 553.71\,s \\
 & \texttt{unif-kkm-1x} &  100.0\% & 10{,}380 & 0.748 & 0.24\,s \\
 & \texttt{unif-kkm-0.5x} & 100.0\% &  5{,}187 & 0.700 & 0.15\,s \\
\bottomrule
\end{tabular}
\label{tab:exp3-summary}
\end{table}

\begin{figure}[h]
\centering
\includegraphics[width=\textwidth]{../results/figures/exp3_fig1_ari_bars.png}
\caption{Mean ARI ($\pm$ std), 3 versions $\times$ 6 algorithms, 10 reps each.}
\label{fig:exp3-fig1}
\end{figure}

\begin{figure}[h]
\centering
\includegraphics[width=\textwidth]{../results/figures/exp3_fig2_tau_bars.png}
\caption{Mean $\tau$ ($\pm$ std, log scale) per algorithm.}
\label{fig:exp3-fig2}
\end{figure}

\begin{figure}[h]
\centering
\includegraphics[width=\textwidth]{../results/figures/exp3_fig3_wallclock_bars.png}
\caption{Mean wall-clock/rep in seconds ($\pm$ std, log scale) per algorithm.}
\label{fig:exp3-fig3}
\end{figure}

\begin{figure}[h]
\centering
\includegraphics[width=\textwidth]{../results/figures/exp3_fig4_convergence_bars.png}
\caption{Convergence rate ($\pm$ std) per algorithm.}
\label{fig:exp3-fig4}
\end{figure}

\paragraph{\texttt{exp3\_1}/\texttt{exp3\_2} confirm KABC's advantage;
\texttt{exp3\_3} is markedly harder --- all at the full 50 reps now.}
\texttt{kabc} reaches ARI $1.000$ on \texttt{exp3\_1} and $0.995$ on
\texttt{exp3\_2} ($98\%$ stopped) while \texttt{acb} is at or near $0.000$
(mean-blind, as expected). On \texttt{exp3\_3} --- tight mean gap,
\emph{no} MMD-maximizing cherry-pick --- even \texttt{kabc} drops to ARI
$0.791$ ($80\%$ stopped): removing the cherry-pick and keeping the mean gap
tight produces genuinely harder sampled pairs for every algorithm,
consistent with the design's intent (an ablation of \emph{why}
\texttt{exp3\_1} works, not just \emph{that} it works). These numbers are
close to the 10-rep pilot's ($1.000$, $1.000$, $0.700$ respectively) ---
the pilot's qualitative picture holds at 50 reps.

\paragraph{TaS-FW is correctness-competitive at a real, large compute
cost, confirmed again at the full 50 reps.} \texttt{tas-fw} reaches ARI
$1.000$ on all three versions, but at a mean wall-clock cost of
409--576\,s/rep --- roughly three orders of magnitude slower than every
other algorithm (all under 1.6\,s/rep). Its own formal stopping rule
rarely fires within budget ($4$--$70\%$ ``stopped'' across versions) even
though its partition estimate is already correct, the same pattern
documented in the prior (superseded) report. \texttt{exp3\_3}'s stopped
rate fell further at 50 reps ($10\%\to4\%$) than the other two versions
moved, consistent with it being the hardest of the three pairing regimes
for every algorithm, not just \texttt{kabc}.

\paragraph{ACB partially recovers on \texttt{exp3\_2}'s wider mean gap.}
Unlike the uniformly-zero ARI on \texttt{exp3\_1}/\texttt{exp3\_3},
\texttt{acb} reaches mean ARI $0.124$ ($20\%$ stopped) on \texttt{exp3\_2}
(moderate-to-large gap) --- some real mean separation gives it partial
signal, though still far behind \texttt{kabc}. (The pilot's $0.215$/$40\%$
was on the high side of what the full 50 reps show; the direction of the
effect --- partial recovery specifically on the wider-gap version --- is
unchanged.)

\FloatBarrier
\section{Results --- Experiment 3 follow-up: tas-fw low-budget probe and calibration}
\label{sec:results-exp3-tasfw-followup}

Ad hoc, NOT part of the frozen Experiment 3 design above --- kept fully
separate in code, config, and results (\texttt{results/exp3\_1|2|3/results.parquet}
were never touched). Prompted by \S\ref{sec:results-exp3}'s observation
that \texttt{tas-fw} almost never reports \texttt{stopped=True} yet its
point estimate is usually already correct.

\paragraph{Stage 1: low-budget probe
(\texttt{scripts/run\_experiment3\_tasfw\_lowbudget.py},
\texttt{results/exp3\_tasfw\_lowbudget/results.parquet}).} Two FIXED,
reduced budgets --- \texttt{tas-fw-0.1x} (\texttt{budget\_cap=1500}) and
\texttt{tas-fw-0.2x} (\texttt{budget\_cap=3000}), i.e.\ 0.1x/0.2x of
production's \texttt{budget\_cap=15000} --- across all 50 reps $\times$ 3
versions, same real movie pairs as production (same \texttt{exp\_id}/
\texttt{seed\_multiplier}).

\begin{table}[h]
\centering
\small
\begin{tabular}{lrr}
\toprule
Version & \texttt{tas-fw-0.1x} mean ARI & \texttt{tas-fw-0.2x} mean ARI \\
\midrule
\texttt{exp3\_1} & $1.000 \pm 0.000$ & $1.000 \pm 0.000$ \\
\texttt{exp3\_2} & $0.990 \pm 0.071$ & $1.000 \pm 0.000$ \\
\texttt{exp3\_3} & $0.526 \pm 0.438$ & $0.776 \pm 0.365$ \\
\bottomrule
\end{tabular}
\caption{Low-budget tas-fw probe, full 50-rep run. \texttt{frac\_stopped=0}
throughout (never certifies even at these tiny budgets).}
\end{table}

\texttt{exp3\_1}/\texttt{exp3\_2}'s point estimate is already essentially
correct at just 10\% of production's budget --- the remaining 90\% is
spent entirely chasing certification, not accuracy. \texttt{exp3\_3} is
genuinely different: real, substantial ARI improvement from more budget
($0.526\to0.776$ doubling $1500\to3000$) --- confirms tas-fw works, and
benefits from more samples, at budgets far smaller than production's
15{,}000, specifically on the harder, non-cherry-picked pairing regime.

\paragraph{Stage 2: calibration (\texttt{scripts/calibrate\_tas\_fw.py}).}
\texttt{tas\_fw.py}'s $\beta(t,\delta)$ (Eq.~66) had no calibratable
constant at all before this. Added one (\texttt{const}, multiplying the
whole bound, same convention as kabc's) via a private
\texttt{\_fit(..., const=1.0)} behind the unchanged public \texttt{fit()}.

Needed its own protocol, not a fourth call into
\texttt{scripts/calibrate\_constants.py}'s live-rerun bisection, for two
structural reasons:

\begin{itemize}
\item \textbf{Calibration instance}: tas-fw is finite-alphabet only (can't
use Exp1's \texttt{GaussianMixtureEnv}), and Experiment 3 has no fixed
reference instance. Solved with a new synthetic environment,
\texttt{kabc.environments.DiscretePMFPairEnv} --- explicit probability
vectors over the real rating support $\{1,\ldots,5\}$, $N{=}8, K{=}2$
matching Exp3's real deployment scale exactly, NOT real MovieLens data.
Three regimes with real, hand-verified KL-gaps spanning a genuine range:
easy$=1.484$, moderate$=0.804$, hard$=0.041$ (mean-gaps $2.3/1.6/0.4$,
straddling Exp3's own $0.25/0.5$ close/far boundaries).
\item \textbf{Cost}: tas-fw's real per-run cost at $\text{budget\_cap}=15{,}000$
is $\sim$400--900s; a literal live-rerun bisection would have cost many
hours. Avoided by exploiting a structural fact specific to tas-fw: neither
the C-tracking sampling rule nor $Z_t$ itself ever reads $\beta$/\texttt{const}
anywhere --- the sampled trajectory is IDENTICAL regardless of
\texttt{const} (unlike kabc, whose threshold gates its own round-based
sampling schedule). So each calibration rep needs exactly ONE real run,
with early exit disabled and the full trajectory $(t, Z_t, \beta_{\text{raw}},
\text{partition})$ logged --- after which ANY candidate \texttt{const}'s
counterfactual stopping time/correctness is recoverable OFFLINE, for free:
$\tau(\text{const}) = \min\{t : Z_t(t) \ge \text{const}\cdot\beta_{\text{raw}}(t)\}$.
\end{itemize}

Ground truth is used ONLY inside this one-time offline script (exactly
like kabc/kabc-sg/acb's own calibration) --- never inside \texttt{tas\_fw.py}'s
\texttt{\_fit}/\texttt{fit} themselves --- and the PMFs are synthetic, so
no real Exp3 data is ever seen during calibration either.

Protocol: 30 reps/regime $\times$ 3 regimes $=$ 90 real trajectory
recordings for the search block, another 90 for held-out validation,
$\delta=0.05$, $\text{budget\_cap}=15{,}000$. A candidate \texttt{const}
passes only if empirical error (not stopped OR wrong) stays
$\le\delta/2$ in EVERY regime independently, not pooled. Real cost: 90+90
runs, wall-clock $5998.4\text{s} + 5995.7\text{s} = 12{,}015.6\text{s}$
total ($\sim$3.34h) --- mean $703.3$s/run, notably higher than an isolated
single-run timing check ($319.3$s), attributed to CPU contention across
the 10 concurrent workers.

\textbf{Result: $\text{const}=0.240940$.} Search-set empirical error was
$0.0$ in all three regimes. \textbf{Held-out validation flagged a real
generalization concern}: the ``moderate'' regime's held-out error was
$0.10$ --- 4$\times$ the required $\delta/2=0.025$ bar --- even though its
search-set error was $0$ (\texttt{held\_out\_generalizes: false} in
\texttt{configs/calibration\_tas\_fw.yaml}). Reported and carried forward
as an open caveat, not silently accepted.

\paragraph{Stage 3: calibrated production run
(\texttt{scripts/run\_experiment3\_tasfw\_calibrated.py},
\texttt{results/exp3\_tasfw\_calibrated/results.parquet}).}
\texttt{tas-fw-calibrated} at Exp3's REAL $\text{budget\_cap}=15{,}000$,
same real movie pairs as production, full 50 reps $\times$ 3 versions.
Despite the held-out caveat above, real deployment showed no sign of it
manifesting as a broken guarantee --- \textbf{zero delta-guarantee
violations across all 150 rows} (every \texttt{stopped=True} run has ARI
exactly $1.0$).

\begin{table}[h]
\centering
\small
\begin{tabular}{lrrr}
\toprule
Version & Uncalibrated (ARI / stopped / mean $\tau$) & Calibrated (ARI / stopped / mean $\tau$) \\
\midrule
\texttt{exp3\_1} & $1.000$ / $38\%$ / $13{,}660$ & $1.000$ / $\mathbf{100\%}$ / $3{,}092$ \\
\texttt{exp3\_2} & $1.000$ / $70\%$ / $9{,}942$ & $1.000$ / $\mathbf{100\%}$ / $2{,}196$ \\
\texttt{exp3\_3} & $1.000$ / $4\%$ / $14{,}928$ & $\mathbf{0.967}$ / $\mathbf{50\%}$ / $11{,}562$ \\
\bottomrule
\end{tabular}
\caption{Calibrated vs.\ uncalibrated tas-fw, real Exp3 data, full 50-rep run.}
\end{table}

\textbf{Honest finding, not a free lunch}: convergence/certification rate
improves dramatically everywhere (4--12.5$\times$), with ARI unchanged on
\texttt{exp3\_1}/\texttt{exp3\_2} but a real, small accuracy cost on
\texttt{exp3\_3} ($1.000\to0.967$, i.e.\ $-3.3$pp) traded for that
version's 12.5$\times$ convergence-rate gain ($4\%\to50\%$) --- a
favorable trade, not a strictly free improvement.

\begin{figure}[h]
\centering
\includegraphics[width=0.85\textwidth]{../results/figures/exp3_fig8_focused_convergence_bars.png}
\caption{Convergence rate, \texttt{kabc}/\texttt{acb}/\texttt{tas-fw}/\texttt{tas-fw-calibrated},
all three versions, full 50-rep run.}
\label{fig:exp3-tasfw-calibrated-convergence}
\end{figure}

\begin{figure}[h]
\centering
\includegraphics[width=0.85\textwidth]{../results/figures/exp3_fig5_focused_ari_bars.png}
\caption{ARI, same four algorithms and versions.}
\label{fig:exp3-tasfw-calibrated-ari}
\end{figure}

\FloatBarrier
\section{Results --- Experiment 4 (clean redesign, complete at 50 reps)}
\label{sec:results-exp4}

\textbf{Complete as of this revision.} Experiment 4 was substantially
redesigned on 2026-09-16 (Section~\ref{sec:setup-exp4}) -- two algorithms,
two real-data-resampled settings, a real-verified memory-guard override
raising kabc's honest convergence ceiling to k$\le$8. The full 50-rep
production run under this design finished cleanly (no errors) on
2026-09-17; \texttt{results/exp4\_baseline/results.parquet} and
\texttt{results/exp4\_amplified\_imbalance/results.parquet} each hold 100
rows (50 reps $\times$ 2 algorithms). The earlier design's 10-rep-pilot
results (\texttt{fixed}/\texttt{resampled}, six algorithm rows) remain
superseded and are not reported here.

\paragraph{Headline.} \texttt{kabc}'s mean ARI stays high and degrades only
slightly from \texttt{baseline} to \texttt{amplified\_imbalance} ($0.964
\to 0.913$); \texttt{unif-kkm-1x}'s mean ARI, at the IDENTICAL matched
per-rep budget (that rep's own realized $\tau_{\text{kabc}}$), drops much
further ($0.668 \to 0.489$). The imbalance-robustness gap this experiment
was designed to isolate is real, and widens exactly as intended when the
far pair's size ratio is pushed from 12:2 to 20:2.

\begin{table}[h]
\centering
\small
\begin{tabular}{llrrr}
\toprule
Setting & Algorithm & Mean ARI ($\pm$ std) & Convergence rate & Mean $\tau$ \\
\midrule
\texttt{baseline} (K=7, N=40, cap=400{,}000) & \texttt{kabc} & $0.964 \pm 0.059$ & $66\%$ (33/50) & 285{,}492 \\
\texttt{baseline} & \texttt{unif-kkm-1x} & $0.668 \pm 0.148$ & $100\%$ (by construction) & 285{,}492 (tied) \\
\texttt{amplified\_imbalance} (K=12, N=76, cap=800{,}000) & \texttt{kabc} & $0.913 \pm 0.083$ & $32\%$ (16/50) & 705{,}535 \\
\texttt{amplified\_imbalance} & \texttt{unif-kkm-1x} & $0.489 \pm 0.079$ & $100\%$ (by construction) & 705{,}535 (tied) \\
\bottomrule
\end{tabular}
\caption{Headline results, full 50-rep run. \texttt{unif-kkm-1x} always
reports \texttt{stopped=True} by construction (it is non-adaptive and
always spends its full assigned budget); its ``convergence rate'' is not a
meaningful statistical-certification signal the way \texttt{kabc}'s is.}
\label{tab:exp4-headline}
\end{table}

\paragraph{kabc's stopped-vs-capped ARI split.} The point estimate is
already good well before the run formally certifies -- the same pattern
documented for \texttt{tas-fw}'s stopping rule on Experiment 3's real data
(Section~\ref{sec:results-exp3}). On \texttt{baseline}:
stopped ($33/50$) ARI $=0.972\pm0.057$; capped ($17/50$) ARI
$=0.950\pm0.062$ -- capping costs almost nothing. On
\texttt{amplified\_imbalance}: stopped ($16/50$) ARI $=0.940\pm0.049$;
capped ($34/50$) ARI $=0.900\pm0.093$ -- still only a small gap despite
K=12's much larger majority-capped share.

\begin{table}[h]
\centering
\small
\begin{tabular}{llrrrr}
\toprule
Setting & Algorithm & \texttt{ari\_close\_pair} & \texttt{ari\_far\_pair} & \texttt{misclass\_large} & \texttt{misclass\_small} \\
\midrule
\texttt{baseline} & \texttt{kabc} & $0.927 \pm 0.237$ & $0.990 \pm 0.039$ & $0.010 \pm 0.039$ & $0.116 \pm 0.154$ \\
\texttt{baseline} & \texttt{unif-kkm-1x} & $0.568 \pm 0.401$ & $0.526 \pm 0.267$ & $0.216 \pm 0.113$ & $0.249 \pm 0.200$ \\
\texttt{amplified\_imbalance} & \texttt{kabc} & $0.903 \pm 0.248$ & $0.994 \pm 0.034$ & $0.052 \pm 0.065$ & $0.223 \pm 0.198$ \\
\texttt{amplified\_imbalance} & \texttt{unif-kkm-1x} & $0.482 \pm 0.339$ & $0.201 \pm 0.076$ & $0.362 \pm 0.066$ & $0.427 \pm 0.147$ \\
\bottomrule
\end{tabular}
\caption{Per-pair ARI (close vs.\ far) and per-size misclassification rate,
full 50-rep run.}
\label{tab:exp4-breakdown}
\end{table}

\paragraph{Key finding: the size-equalization bias, isolated.}
\texttt{unif-kkm-1x}'s \texttt{ari\_far\_pair} COLLAPSES under amplified
imbalance ($0.526 \to 0.201$), even though the far pair is BY CONSTRUCTION
the easier-separated of the two ($|\text{mean}_A-\text{mean}_B|\ge 0.5$,
vs.\ the close pair's $\le 0.25$) -- direct evidence that k-means'
size-equalization bias, not separation difficulty, drives the failure:
\texttt{misclass\_rate\_large} nearly doubles ($0.216 \to 0.362$) as the
size ratio worsens from 12:2 to 20:2. \texttt{kabc} shows the opposite
pattern -- \texttt{ari\_far\_pair} if anything IMPROVES slightly ($0.990
\to 0.994$) and \texttt{misclass\_rate\_large} stays small ($0.010 \to
0.052$) -- its decision rule exhibits no comparable size-driven failure
mode. \texttt{kabc}'s own convergence rate (fraction of reps reaching a
fully CERTIFIED stop, not point-estimate quality) drops as expected with K
($66\% \to 32\%$, since \texttt{amplified\_imbalance}'s K=12/N=76 needs
more of the doubling schedule and hits \texttt{budget\_cap} before
finishing more often) -- but as the stopped-vs-capped split above shows,
this does not translate into a comparable ARI cost.

\begin{figure}[h]
\centering
\includegraphics[width=0.85\textwidth]{../results/figures/exp4_fig1_ari_bars.png}
\caption{Mean ARI ($\pm$ std), \texttt{kabc} vs.\ \texttt{unif-kkm-1x}, both
settings, full 50-rep run.}
\label{fig:exp4-ari}
\end{figure}

\begin{figure}[h]
\centering
\begin{subfigure}{0.48\textwidth}
\includegraphics[width=\textwidth]{../results/figures/exp4_fig0b_instance_overview_baseline.png}
\caption{\texttt{baseline} (K=7, N=40).}
\end{subfigure}
\hfill
\begin{subfigure}{0.48\textwidth}
\includegraphics[width=\textwidth]{../results/figures/exp4_fig0b_instance_overview_amplified_imbalance.png}
\caption{\texttt{amplified\_imbalance} (K=12, N=76).}
\end{subfigure}
\caption{Instance overview at rep=0 (arm counts per movie, pairwise
mean-separation heatmap, pairwise MMD$^2$ heatmap), both settings.}
\label{fig:exp4-overview}
\end{figure}

\FloatBarrier
\section{Inferences and Conclusions}
\label{sec:conclusions}

\paragraph{6.1 KABC's empirical error remains exactly zero on Experiments 1
and 2, at every individual grid point/setting, not just in aggregate ---
all three experiments now at the full 50 reps; Experiment 3's hardest
version (\texttt{exp3\_3}) shows KABC is not universally immune.}
Across all 14 Experiment~1 grid points and all 6 Experiment~2 settings,
\texttt{kabc}'s error rate is $0.0\%$ and its ARI is $1.000$ --- checked
per-point/per-setting. On Experiment~3
(now at the full 50 reps), \texttt{kabc} is likewise near-perfect on
\texttt{exp3\_1} ($1.000$) and \texttt{exp3\_2} ($0.995$, $98\%$ stopped),
but drops to $79.1\%$ ARI on \texttt{exp3\_3} --- the version specifically
designed to remove the MMD-maximizing cherry-pick while keeping the mean
gap tight, essentially unchanged from the pilot's $70\%$. This is a
genuinely useful finding: KABC's near-universal correctness on synthetic
data does not automatically transfer to every real-data regime;
\texttt{exp3\_3} is where that boundary shows up.

\paragraph{6.2 KABC-SG's fragility is confirmed again under the redesigned
grids, and is clearly family-/point-dependent rather than uniform --- all
three experiments now at the full 50 reps.} In Experiment~1, \texttt{kabc-sg}
fails to converge at the two hardest $d$ points ($d\in\{50,100\}$: $0/50$
reps stopped at either, $\tau$ pinned exactly at the doubling trick's
$277{,}440$ ceiling) and converges only $8\%$ of the time at $\Delta=1$,
while converging easily elsewhere. In Experiment~2, the pattern is
different but equally sharp: \texttt{kabc-sg} converges \emph{often} on
\texttt{moment\_matched\_mixture} but its ARI there is markedly lower than
its convergence rate would suggest --- a confident-wrong-answer failure
mode, distinct from Experiment~1's budget-exhaustion failure mode on the
same algorithm. Precisely: $83.2\%$ ($79/95$) of \texttt{kabc-sg}'s
\texttt{stopped=True} runs on \texttt{moment\_matched\_mixture} have
$\text{ARI}\neq 1$, a direct violation of its calibrated $\delta=0.05$
guarantee by over an order of magnitude, while \texttt{kabc} shows zero
errors on the identical family/budget (Section~\ref{sec:results-exp2}
discusses this in full). The 50-rep aggregate ($83.0\%$ stopped, ARI
$0.680$) is consistent with the pilot's $86.7\%$/$0.616$ --- the
confident-wrong-answer pattern is not a pilot-scale artifact.

\paragraph{6.3 ACB is not merely imprecise on same-mean data --- it is
exactly, provably blind on Experiment 2, and largely (not entirely) blind
on Experiment 3's real pairs. Both now confirmed at the full 50 reps.}
\texttt{acb}/\texttt{acb-0.5x}'s ARI is exactly $0.000$ and stopped rate
exactly $0\%$ at every one of Experiment 2's 6 settings, with $\tau$
landing exactly on the theoretical \texttt{budget\_cap} bound each time
--- unchanged from the pilot, now over 50 reps/setting instead of 10. On
Experiment 3, \texttt{acb} is similarly $0.000$ on
\texttt{exp3\_1}/\texttt{exp3\_3} (tight-gap versions) but reaches a
partial $0.124$ ARI ($20\%$ stopped) on \texttt{exp3\_2} (wider gap) ---
consistent with the theoretical picture: ACB is blind specifically to
same-mean structure, not to clustering in general, and recovers some
signal once real mean separation exists.

\paragraph{6.4 The \texttt{-0.5x} self-comparison rows show real,
non-trivial margin for every algorithm tested, in all three experiments
--- and the rename from 0.9x fixed a real methodological problem, not just
a cosmetic one --- all three experiments now at the full 50 reps.}
\texttt{kabc}'s own ARI at 50\% of its realized budget drops to $0.047$
(Experiment~1), $0.123$ (Experiment~2), and $0.14$ across all three of
Experiment 3's versions --- far from each setting's $1.000$ (or
near-$1.000$) at full budget. \texttt{kabc-sg} and \texttt{acb} show the
same direction of
effect. Because KABC's doubling-trick schedule concentrates roughly half
of its total realized budget in the final round, the earlier 0.9x
multiplier could not reliably distinguish ``genuinely less budget'' from
``one round short of the same checkpoint a much larger cut would also
miss'' --- 0.5x resolves this ambiguity and gives a consistent,
interpretable margin signal across every algorithm and every experiment.

\paragraph{6.5 The most striking new finding this round: \texttt{unif-kkm}
barely degrades under the same 0.5x cut that collapses \texttt{kabc}, and
this reflects a real, informative gap between point-estimate accuracy and
$\delta$-PAC-certified accuracy, not a flaw in the comparison. Confirmed
at both Experiment 1 and 2's full 50 reps.} \texttt{unif-kkm-1x} vs.\
\texttt{unif-kkm-0.5x}: ARI $1.000\to1.000$ (Experiment 1, exact),
$0.987\to0.986$ (Experiment 2, essentially exact) --- while \texttt{kabc}
itself collapses on the identical budget cut ($1.000\to0.047$ and
$1.000\to0.123$ respectively). \texttt{unif-kkm} is a
one-shot heuristic point estimate: once its per-arm sample size is large
enough that embedding noise is small relative to the true cluster
separation, its kernel $k$-means decision saturates and more samples don't
change the discrete answer. \texttt{kabc}, in contrast, does not stop
until it can \emph{formally certify} correctness at confidence level
$\delta$ via its pairwise threshold test --- a much stricter bar, and one
whose doubling-trick schedule means half its own budget usually falls
short of the round that achieves that certification. The practical
reading: \emph{the correct partition is discoverable with far fewer
samples than KABC needs to prove it is correct} --- evidence of how
conservative a rigorous confidence guarantee has to be relative to what a
good heuristic guess needs, not evidence that adaptivity ``doesn't
matter.'' Worth stating explicitly in the paper write-up.

\paragraph{6.6 TaS-FW remains correctness-competitive at a real, now
precisely re-quantified compute cost, confirmed at Experiment 3's full 50
reps.} On all three Experiment~3 versions, \texttt{tas-fw} reaches ARI
$1.000$ (matching or exceeding \texttt{kabc}), at 409--576\,s/rep ---
roughly three orders of magnitude slower than every other algorithm (all
under 1.6\,s/rep). This reproduces the prior (superseded) report's finding
under the current REVISION~4/0.5x design, confirming it isn't an artifact
of the earlier config or of the 10-rep pilot's smaller sample.

\paragraph{6.7 \texttt{acb-gaussian}, added after the main 50-rep run
completed, reaches exactly zero error at a substantially lower budget than
every other algorithm, uniformly across the grid --- the paper's own
``sharper constants'' claim confirmed directly, not just reproduced
qualitatively.} Section~4/Appendix A of \cite{acb2024} describes a
Gaussian-specialized, non-adaptive variant of ACB (known $\Delta$/$\theta$,
every concentration-inequality constant replaced by an exact
$\delta$-level quantile) as the actual procedure behind the paper's own
Figure~1 experiment ($\le 1\%$ error at $\delta=0.1$). Implemented here as
\texttt{acb-gaussian} (Section~\ref{sec:acb-gaussian}) and run on all 700
(grid point, rep) combinations of Experiment~1: $100\%$ stopped, ARI
$1.000$ \emph{exactly}, on every single run --- zero error, not $\le 1\%$.
Its budget advantage over \texttt{kabc} is genuinely uniform across the
grid ($5.4\times$--$17.4\times$ per grid point, aggregate $8.4\times$
squarely inside that range); the advantage over \texttt{kabc-sg}
($4.7\times$--$343\times$ per grid point, aggregate $69\times$) is real
but concentrated at $d\in\{50,100\}$, where \texttt{kabc-sg} also fails to
converge (\S6.2) --- see \S\ref{sec:results-exp1}'s discussion for the
full per-point breakdown and why a single aggregate ratio would have
repeated the same aggregation mistake the ACB per-$\Delta$ error breakdown
(\S6.3) was added to correct. This is only a fair, informative comparison
because $\Delta$ and $\theta$ genuinely are known at Experiment~1's
instance-construction time (Section~\ref{sec:acb-gaussian} states this
explicitly, not as an implicit privilege) --- but taken on those terms, it
is direct empirical evidence that the accuracy/budget frontier this
benchmark's other algorithms sit on is set by their conservative,
proof-derived thresholds, not by any fundamental difficulty of the
underlying instances at this scale.

This result also completes a before/after picture worth stating
explicitly rather than leaving implicit. \texttt{acb}'s architecture
(SRI+ADC, \S\ref{sec:baselines}) \emph{does} have a graceful fallback
mechanistically identical to \texttt{kabc}'s own
(\S\ref{sec:robustness}): when SRI fails to find $K$ representatives
within budget, \texttt{acb} returns the trivial one-cluster partition with
\texttt{stopped=False}, scoring ARI exactly $0.000$ by the same
construction. This is why Experiment~2's \texttt{acb}/\texttt{acb-0.5x}
rows (\S\ref{sec:results-exp2}) show exactly $0.000$ ARI at $0\%$ stopped:
$\Delta_*=0$ exactly by that experiment's own same-mean construction, so
SRI \emph{provably} cannot succeed --- an honest non-convergence, not a
miscalibration symptom, and not the mechanism this paragraph is about. The
real silent-error failure mode sits one stage later: once SRI \emph{does}
succeed (\texttt{stopped=True} territory), ADC's classification depends
on its own sample-size constants ($I$, $J$, Eq.~13) being large enough for
that instance's actual separation --- and because \texttt{acb}'s single
shared \texttt{c\_cal} multiplier (\S\ref{sec:calibration}) is calibrated
against one fixed instance, not validated at every hard grid point, that
classification can be confidently wrong with no non-convergence signal to
flag it. This is directly evidenced, not merely plausible:
Experiment~1's smallest gaps (\S\ref{sec:results-exp1}) show
\texttt{acb} at $100\%$ \texttt{stopped} throughout, yet error reaching
$18\%$ at $\Delta\in\{1,2\}$ --- SRI succeeded, ADC classified, and was
wrong on nearly one in five runs. Experiment~3's \texttt{exp3\_2}
(\S\ref{sec:results-exp3}) shows the identical pattern on real data: among
its 10 \texttt{stopped=True} reps, 5 have ARI $\ne 1$ (values $-0.12$,
$0.16$ twice, $0.495$ twice) against 5 exactly correct, while all 40
non-stopped reps score exactly $0.000$ --- the same
honest-fallback/silent-error split, on the same algorithm, on real data.
\texttt{acb-gaussian} is the same SRI+ADC architecture with its
classification-stage constants sized exactly right instead of borrowed
generically: identical decision structure, zero silent errors across all
700 runs, and the cheapest budget in the whole benchmark. The lesson is
not that ACB's architecture is unreliable --- its own honest-fallback path
works exactly as designed --- it is that this specific instantiation's
classification-stage constants were never calibrated for non-Gaussian,
non-oracle-parameter settings, a point worth carrying into the discussion
with \cite{acb2024}'s authors alongside Experiment~2/3's own \texttt{acb}
findings.

\paragraph{6.8 All three experiments are now complete at the full planned
50 reps; the pilot's qualitative conclusions held up throughout.}
Experiments 1, 2, and 3 are all at 50 reps/algorithm/grid-point (6{,}300 +
2{,}400 + 900 = 9{,}600 total rows across the three experiments, including
\texttt{acb-gaussian}'s 700 rows merged into Experiment~1 after its main
run completed --- \S\ref{sec:acb-gaussian}), and every observation above
(\S6.1--6.7) reflects that complete data. Where a
10-rep-pilot number is quoted alongside its 50-rep counterpart in the
sections above, the two are consistently close --- no pilot-scale finding
in this report reversed direction once the full data landed; the two
genuine surprises the pilot's small sample size produced (Experiment~1's
apparently ``well-behaved'' \texttt{kabc-sg} at $N=90$, and the exact
pilot-vs-50-rep ACB/\texttt{kabc-sg} percentages) were both explicitly
re-checked against the full data and resolved (\S\ref{sec:results-exp2}
above; the $N=90$ finding is confirmed as pilot-scale noise, not a real
effect --- see the Experiment~1 results section above). One caveat
remains open regardless of rep count: \texttt{kabc-sg}'s and
\texttt{acb}'s calibrated constants (Section~\ref{sec:calibration}) were
validated against a single fixed calibration instance and are not
re-validated against every hard point in these three experiments --- the
error/ARI numbers throughout this report should be read as ``as
calibrated in this repo,'' not as a claim about the underlying
algorithms' theoretical guarantees.

\FloatBarrier
\section{Appendix: why Experiment 4's \texttt{budget\_cap} increase and its memory-guard override are two independent changes}
\label{sec:appendix-exp4-budget-vs-guard}

Both \texttt{baseline}'s $\text{budget\_cap}=400{,}000$ and
\texttt{amplified\_imbalance}'s $\text{budget\_cap}=800{,}000$
(Section~\ref{sec:setup-exp4}) are sometimes read as ``caused by'' the
memory-guard override. That is only half right --- the two are separate,
independently-necessary changes, and neither alone would have raised
kabc's convergence ceiling from k$\le$7 to k$\le$8.

\paragraph{The \texttt{budget\_cap} value itself comes from kabc's own
doubling-trick cost schedule, not from the memory-guard formula.} For a
given $N$, the doubling trick's cumulative pull cost through iteration $k$
is fixed by \texttt{thresholds.py}'s \texttt{iteration\_budget}/
\texttt{iteration\_confidence} (Algorithm~2), independent of any
memory-guard parameter. At $N{=}40$ (\texttt{baseline}), the cumulative
cost through k=8 is $360{,}520$; at $N{=}76$
(\texttt{amplified\_imbalance}), it is $735{,}148$. Each setting's
\texttt{budget\_cap} was set comfortably above that number ($400{,}000$:
$\approx$11\% margin; $800{,}000$: $\approx$8.8\% margin) specifically so
that \texttt{budget\_cap} itself would no longer be the thing stopping
kabc at k=7 --- exactly the way the earlier design's
$\text{budget\_cap}=200{,}000$ did (\S\ref{sec:setup-exp4}'s superseded
version): $200{,}000 \ll 360{,}520$, so k=8 was financially unaffordable
under the old cap regardless of memory-guard settings.

\paragraph{The memory-guard override answers a categorically different
question: is it even SAFE to attempt k=8's Gram-matrix computation at
all.} Under the original guard (11 workers, 4$\times$ safety, trip point
$n_k>2{,}852$), k=8's per-arm sample size ($n_k{\approx}4{,}602$ at
$N{=}40$; $n_k{\approx}4{,}933$ at $N{=}76$) would have been refused
outright (\texttt{\_MemoryGuardExceeded}) regardless of how large
\texttt{budget\_cap} was set. The real-verified override
(\texttt{memory\_guard\_n\_workers=5},
\texttt{memory\_guard\_safety\_factor=2}, both matching this experiment's
actual $n_{\text{jobs}}{=}5$ and its real measured peak RSS --- see
Section~\ref{sec:setup-exp4}) raised that trip point to $n_k>5{,}982.6$,
which is what makes attempting k=8 safe in the first place.

\paragraph{Both changes were necessary; neither alone was sufficient.}
Old \texttt{budget\_cap} ($200{,}000$) with the new guard: still blocked
at k=7 (can't afford k=8's cost). Old guard ($n_k>2{,}852$) with the new,
larger \texttt{budget\_cap}: still blocked at k=7 (the guard refuses to
even attempt k=8's Gram matrices, before any budget check on that
iteration is reached). Only raising both together unlocked k=8 for both
settings.

\section{Appendix: how loose is Experiment 1's \texttt{budget\_cap} relative to the paper's own Theorem 3.1 bound}
\label{sec:appendix-exp1-budgetcap-vs-eq8}

Experiment 1's \texttt{budget\_cap} (\texttt{s\_star\_formula},
\texttt{constant=200}: $\texttt{budget\_cap} = 200 \cdot N \cdot s_*^{-2}
\cdot \log(N/\delta)$, \texttt{scripts/common.py}'s
\texttt{\_resolve\_budget\_cap}) is a practical, hand-chosen safety
multiplier, not a value derived from the paper's own proof. This note
quantifies exactly how it compares to the paper's fully-derived, proven
one-shot sample-complexity bound (Theorem 3.1, Eq.~8):
$$8N\left(\frac{128}{s_*^2}\vee 1\right)\log\!\left(\frac{32(N^2-N)k_*^2}{\delta}\right),
\qquad k_* = \left\lceil \log_2\!\frac{128}{s_*^2}\right\rceil \vee 1,$$
using the \emph{same} Monte-Carlo $s_*^2$ estimate
(\texttt{kabc.metrics.estimate\_s\_star\_sq}, base config, fixed seed) for
both formulas, so the comparison isolates the constant/log-term
difference rather than a different difficulty measure. Computed for
Experiment 1's base config ($N{=}60$, $K{=}3$, $d{=}10$, $\sigma{=}1$,
$\delta{=}0.05$) across its main $\Delta$ sweep:

\begin{center}
\begin{tabular}{rrrrrr}
\toprule
$\Delta$ & $s_*^2$ (MC est.) & \texttt{budget\_cap} (ours) & $k_*$ & Eq.~8 bound (paper) & ratio (ours / Eq.~8) \\
\midrule
1   & 0.146894 & 579{,}200 & 10 & 8{,}046{,}726 & 0.071980 \\
1.5 & 0.329325 & 258{,}350 & 9  & 3{,}549{,}893 & 0.072777 \\
2   & 0.582866 & 145{,}971 & 8  & 1{,}980{,}894 & 0.073689 \\
3   & 1.185555 & 71{,}765  & 7  & 960{,}046     & 0.074752 \\
4   & 1.401399 & 60{,}712  & 7  & 812{,}179     & 0.074752 \\
6   & 1.613719 & 52{,}724  & 7  & 705{,}319     & 0.074752 \\
\bottomrule
\end{tabular}
\end{center}

\paragraph{Finding.} Our \texttt{budget\_cap} is consistently
$\approx$7.2--7.5\% of the paper's exact Eq.~8 bound across the entire
$\Delta$ sweep (i.e.\ Eq.~8 is $\approx$13.4--13.9$\times$ larger), and the
ratio is remarkably stable, converging as $\Delta$ grows and $k_*$
settles at 7.

\paragraph{This is a distinct knob from \texttt{c\_cal}
(Section~\ref{sec:setup-exp4}'s memory-guard note and the main text's
calibration protocol), not the same looseness measured twice.}
\texttt{budget\_cap} is a hard spending \emph{ceiling} passed to
\texttt{\_fit} --- it bounds how much \texttt{tau} an algorithm is
\emph{allowed} to spend before being forcibly stopped, and is rarely
binding for kabc in practice (kabc's own calibrated stopping rule usually
halts well under cap). \texttt{c\_cal} ($\approx 0.103516$ for kabc, see
\texttt{configs/calibration.yaml}) instead multiplies the per-iteration
same/different-cluster decision bound (Eq.~6) inside
\texttt{\_cluster\_once} --- it governs \emph{when kabc actually decides
to stop}. The two ratios ($\approx$0.072--0.075 for
\texttt{budget\_cap}/Eq.~8, $\approx$0.103 for \texttt{c\_cal}) land in a
similar order of magnitude by coincidence, not because they are the same
quantity measured two ways.

\begin{thebibliography}{9}

\bibitem{kabc2026}
Thuot, Vogt, Ghoshdastidar, Verzelen.
\textit{Kernel-based Active Clustering of Bandit Arms}.
arXiv:2601.07535, 2026.

\bibitem{acb2024}
Thuot, Carpentier, Giraud, Verzelen.
\textit{Active clustering with bandit feedback}.
arXiv:2406.11485, 2024.

\bibitem{tasfw2024}
Yavas, Huang, Tan, Scarlett.
\textit{A General Framework for Clustering and Distribution Matching with Bandit Feedback}.
arXiv:2409.05072, 2024.

\bibitem{wolferalquier2025}
Wolfer, Alquier.
\textit{Variance-aware estimation of kernel mean embedding}.
Journal of Machine Learning Research, 26(57):1--48, 2025.

\bibitem{tolstikhin2017}
Tolstikhin, Sriperumbudur, Muandet.
\textit{Minimax estimation of kernel mean embeddings}.
Journal of Machine Learning Research, 18(86):1--47, 2017.

\bibitem{boc2024}
Yang, Zhong, Tan.
\textit{Optimal Clustering with Bandit Feedback}.
Journal of Machine Learning Research, 25(186):1--54, 2024. arXiv:2202.04294.

\bibitem{banditclustering2025}
Graf, Thuot, Verzelen.
\textit{Clustering Items through Bandit Feedback: Finding the Right Feature out of Many}.
Proceedings of the 42nd International Conference on Machine Learning (ICML), 2025. arXiv:2503.11209.

\bibitem{clusteredbai2025}
Yash, Karamchandani, Ghosh.
\textit{Near Optimal Best Arm Identification for Clustered Bandits}.
Proceedings of the 42nd International Conference on Machine Learning (ICML), 2025. arXiv:2505.10147.

\bibitem{atboc2025}
Chandran, G. Dhinesh; Reddy, Kota Srinivas; Bhashyam, Srikrishna.
\textit{Online Clustering of Data Sequences with Bandit Information}.
2025 IEEE International Symposium on Information Theory (ISIT), 2025. arXiv:2501.11421.

\bibitem{ebc2026}
Chandran, G. Dhinesh; Reddy, Kota Srinivas; Bhashyam, Srikrishna.
\textit{Efficient Clustering in Stochastic Bandits}.
arXiv:2601.09162, 2026 (preprint).

\end{thebibliography}

\end{document}
